\chapter{Einleitung}
\label{chap:intro}
Reinforcement Learning (RL) hat sich in den letzten Jahren von theoretischen Konzepten zu praxisnahen Anwendungen entwickelt. Autonome Roboter, die eigenständig Hindernisse umfahren, und Algorithmen im Finanzhandel, die Markttrends selbstständig identifizieren, sind praxisnahe Beispiele für den Erfolg von RL. Zentrales Merkmal dieser Methode ist das eigenständige Erlernen von Strategien durch Agenten aus Interaktionen mit ihrer Umgebung, ohne dass ein vollständiges Modell der Umgebung vorab benötigt wird. Trotz dieser beeindruckenden Fortschritte gestaltet sich der Lernprozess häufig komplex und instabil, da die Trainingsdaten nicht stationär sind, sondern unmittelbar von der sich stetig ändernden Policy des Agenten abhängen \cite{sutton2018reinforcement}.

Während viele erfolgreiche Anwendungen auf Single-Agent-Setups beruhen, bietet Multi-Agent Reinforcement Learning (MARL) die Möglichkeit, komplexere Szenarien mit mehreren kooperativen oder konkurrierenden Agenten abzubilden. MARL hat somit großes Potenzial für die Untersuchung emergenter Phänomene, wie etwa Teamstrategien oder komplexer Interaktionsdynamiken. Die gleichzeitige Optimierung mehrerer Policies führt jedoch zu zusätzlichen Herausforderungen, insbesondere im Bereich des Credit Assignment \cite{sutton2018reinforcement}, da unklar bleibt, welche Aktionen einzelner Agenten maßgeblich zum Erfolg beitragen.
\section{Beitrag und Ziele}

In einem ersten Versuch wurde ein MARL-Szenario untersucht, in dem sowohl Jäger- als auch Beute-Agenten \footnote{In dieser Arbeit werden die Begriffe Jäger und Beute verwendet, um die Rollen von Predator und Prey zu bezeichnen. Diese Begriffe sind inhaltlich äquivalent und werden zur sprachlichen Konsistenz im gesamten Text genutzt.} in einer kontinuierlichen 2D-Umgebung trainiert wurden. Dabei sollte insbesondere das emergente Verhalten der Agenten analysiert werden. Aufgrund der hohen Komplexität und der starken Interaktionseffekte zwischen den Agenten zeigte sich jedoch keine stabile Lernkonvergenz. Die Ergebnisse verdeutlichten, dass der Einsatz einer Multi-Agent-Umgebung - ohne vorherige isolierte Betrachtung grundlegender - Fähigkeiten zu instabilen Trainingsbedingungen führt. Diese Erfahrungen unterstreichen die Notwendigkeit zunächst spezifische grundlegende Lernmechanismen und Trainingsbedingungen systematisch in vereinfachten Single-Agent-Szenarien zu analysieren.

Ausgehend von diesen Erkenntnissen ist es das Ziel dieser Arbeit, den Einfluss von gezielten Trainingsszenarien systematisch zu untersuchen. Bengio et al. \cite{10.1145/1553374.1553380} haben gezeigt, dass Curriculum Learning, also die schrittweise Erhöhung der Aufgabenschwierigkeit, Lernprozesse beschleunigt und verbessert. Aufbauend auf diesem Ansatz untersucht diese Arbeit systematisch, wie sich ein modularer Trainingsansatz auf einen Reinforcement-Learning-Agenten auswirkt: Die Subtasks Navigation, Exploration und Flucht werden zunächst separat optimiert und anschließend in ein komplexes Survival-Szenario übertragen. Im Mittelpunkt steht die Frage, ob dieses Vortraining Lernstabilität und -effizienz verbessert, und welche Rolle dabei verschiedene Reward-Funktionen spielen. Ebenfalls wird hierbei untersucht, ob angepasste Modelle ihre Performance in den Subtasks beibehalten, oder ob neue Strategien alte Fähigkeiten verdrängen.

Aktuelle Arbeiten, wie die Studie von Baker et al. \cite{baker2020emergenttoolusemultiagent}, zeigen, dass komplexe Fähigkeiten durch Auto-Curricula entstehen können. Jedoch bleibt dabei häufig offen, welche Teilaufgaben konkret zum Lernfortschritt beitragen und wie gut die gelernten Fähigkeiten auf neue Situationen übertragbar sind. Zhuang et al. \cite{zhuang2020comprehensivesurveytransferlearning} betonen zusätzlich, dass Transfer Learning und die Wiederverwendung erlernter Fähigkeiten entscheidend für die Effizienz in komplexen RL-Szenarien sind.

Daher wird in dieser Arbeit untersucht, wie sich das isolierte Vortraining auf spezifische Subtasks auf die Gesamtleistung des Agenten in einem komplexen Überlebensszenario auswirkt. Ziel ist es zu analysieren, inwiefern diese modulare Trainingsstrategie zu stabileren und effizienteren Lernprozessen führt, und ob die erlernten Fähigkeiten sinnvoll auf komplexere Umgebungen übertragen werden können. Hierbei wird insbesondere  die Auswirkung verschiedener Rewardfunktionen auf den Lernprozess analysiert. Ebenfalls wird hierbei untersucht, ob angepasste Modelle ihre Performance in den Subtasks beibehalten.

Ein weiterer Beitrag dieser Arbeit ist die Entwicklung eines öffentlich verfügbaren Benchmark-Frameworks auf Basis der PettingZoo-Bibliothek \cite{Terry_PettingZoo_Gym_for}. Es bietet eine modulare Simulationsumgebung und vorkonfigurierte Trainingsskripte für LSTM-PPO-Agenten und ermöglicht so eine reproduzierbare Grundlage für weiterführende Forschung und komplexere Multi-Agent-Szenarien.

\section{Aufbau der Arbeit}
Im Folgenden werden zunächst die relevanten theoretischen Grundlagen (Kapitel ~\ref{chap:theorie}) vorgestellt. Danach werden die Ergebnisse der Multi-Agenten-Experimente erläutert (Kapitel ~\ref{chap:explorative}). Kapitel ~\ref{chap:methodology} beschreibt die gemeinsame Methodik, während Kapitel ~\ref{chap:single-tasks} die Ergebnisse der Single-Agenten- und Subtask-Experimente darlegt. Abschließend fasst Kapitel ~\ref{chap:fazit} die Resultate zusammen und Kapitel ~\ref{chap:ausblick} skizziert Perspektiven für weiterführende Forschung.